\begin{examplebox}
\par\noindent\begin{minipage}{\linewidth}
\textbf{\textcolor{CExample}{例1：知道三边，求半径（先看能否用勾股定理）}}

\textbf{题目：}三边长为 $5,12,13$ 的三角形，求内切圆半径。
\PIncenterFigure{0}{12}{5}{0.43}{5}{13}{12}{0}{例1：两直角边 $5,12$，斜边 $13$；圆的半径 $r$ 待求。}
\end{minipage}\par
\begin{enumerate}[leftmargin=*,itemsep=0.45em]
 \item $5^2+12^2=13^2$，所以它是直角三角形，两条直角边是 $5,12$。
 \item 面积 $S=\frac12\times5\times12=30$，周长 $L=5+12+13=30$。
 \item $r=\dfrac{2S}{L}=\dfrac{60}{30}=\boxed{2}$。
\end{enumerate}
\textbf{再快一步：}直角三角形专用式
$r=(5+12-13)/2=2$。\textbf{识别信号：}三边满足勾股定理。

\medskip
\par\noindent\begin{minipage}{\linewidth}
\textbf{\textcolor{CExample}{例2：知道半径和周长，求面积（可实现的数据）}}

\textbf{题目：}一个三角形的内切圆半径为 $r=2$，周长为 $L=24$，求面积。
\PIncenterFigure{0}{8}{6}{0.58}{6}{10}{8}{0}{例2的一种可行图形：$L=24$，$r=2$；所求是整个三角形面积。}
\end{minipage}\par
\begin{enumerate}[leftmargin=*,itemsep=0.4em]
 \item 半周长 $p=L/2=12$。
 \item $S=rp=2\times12=\boxed{24}$。
 \item \textbf{验证题目确实可能成立：}例如两直角边为 $6,8$、斜边为 $10$ 的三角形：
 周长 $6+8+10=24$，面积 $6\times8/2=24$，
 半径 $2\times24/24=2$，全部吻合。
\end{enumerate}
\textbf{注意：}半径、周长必须来自真实三角形；数据可行性检查见“易错点”。

\medskip
\par\noindent\begin{minipage}{\linewidth}
\textbf{\textcolor{CExample}{例3：不是直角三角形，三边求半径（海伦公式）}}

\textbf{题目：}三边长为 $13,14,15$ 的三角形，求内切圆半径。
\PIncenterFigure{9}{12}{14}{0.43}{14}{13}{15}{0}{例3：$BC=14,CA=13,AB=15$；先算面积，再求蓝色半径。}
\end{minipage}\par
\begin{enumerate}[leftmargin=*,itemsep=0.45em]
 \item 先算半周长：$p=(13+14+15)/2=21$。
 \item 用海伦公式：
 \[
 \begin{aligned}
 S&=\sqrt{p(p-13)(p-14)(p-15)}\\
  &=\sqrt{21\times8\times7\times6}
   =\sqrt{7056}=84.
 \end{aligned}
 \]
 \item 因而 $r=S/p=84/21=\boxed{4}$。
\end{enumerate}
\textbf{路线判断：}不是直角三角形时，不要硬猜高；三边已知可用海伦公式。

\medskip
\par\noindent\begin{minipage}{\linewidth}
\textbf{\textcolor{CExample}{例3变式：同样的三边，直接求切点分段}}

\textbf{题目：}仍取 $BC=14,CA=13,AB=15$，内切圆切 $BC,CA,AB$
于 $D,E,F$。求 $AE,AF,BF,BD,CD,CE$。
\PTangentFigure{7}{8}{6}
\end{minipage}\par
半周长为 $p=21$，每个顶点的切线段长等于 $p$ 减去其对边长：
\[
\begin{aligned}
 AE=AF&=p-BC=21-14=\boxed{7},\\
 BF=BD&=p-CA=21-13=\boxed{8},\\
 CD=CE&=p-AB=21-15=\boxed{6}.
\end{aligned}
\]
\textbf{回验：}$BD+DC=8+6=14$，$AE+EC=7+6=13$，
$AF+FB=7+8=15$，恰好还原三边。
\textbf{与上一问的区别：}求这些沿边的长度，不需要先用海伦公式求面积，也不用计算半径。

\medskip
\par\noindent\begin{minipage}{\linewidth}
\textbf{\textcolor{CExample}{例4：不求半径，直接求面积比}}

\textbf{题目：}在 $\triangle ABC$ 中，$BC=6,CA=8,AB=10$，
内心为 $I$。求
$S_{\triangle IBC}:S_{\triangle ICA}:S_{\triangle IAB}$。
\PIncenterFigure{6}{8}{6}{0.62}{6}{8}{10}{1}{例4：蓝、橙、绿三块的高同为 $r$；面积比直接读对应底边比。}
\end{minipage}\par

\begin{enumerate}[leftmargin=*,itemsep=0.4em]
 \item 连接 $IA,IB,IC$ 后，三个小三角形到底边的高都是内切圆半径 $r$。
 \item 因而面积比分别为
 \[
  \frac{6r}{2}:\frac{8r}{2}:\frac{10r}{2}
  =6:8:10=\boxed{3:4:5}.
 \]
\end{enumerate}
\textbf{为何更快：}题目只问比值，$r$ 相同会约去，不需要计算半径。

\medskip
\par\noindent\begin{minipage}{\linewidth}
\textbf{\textcolor{CExample}{例5：反过来用面积比找边长}}

\textbf{题目：}在 $\triangle ABC$ 中，$I$ 是内心，已知
$S_{\triangle IBC}:S_{\triangle ICA}=3:5$，并且 $BC=9$，求 $CA$。
\PIncenterFigure{0}{12}{9}{0.43}{9}{15}{12}{1}{例5的一种可行形状（$AB$ 不是题目给定量）：\\[3pt]蓝块与橙块同高，面积比 $3:5$ 对应底边比 $BC:CA=3:5$。}
\end{minipage}\par

\begin{enumerate}[leftmargin=*,itemsep=0.4em]
 \item 因两块小三角形的高都是 $r$，
 $S_{\triangle IBC}:S_{\triangle ICA}=BC:CA=3:5$。
 \item 代入 $BC=9$，得到 $9:CA=3:5$，因此
 \[
  CA=\frac{9\times5}{3}=\boxed{15}.
 \]
\end{enumerate}
\textbf{核心作用：}把面积比变成边长比，不必求出两个面积。
\end{examplebox}
